% GDZ PPN236018647,native595lower–597;HKUST596–598.
\NumberedItem{134}
\noindent\hspace*{\parindent}Demonstrationem omnino nova[m] theorematis fundamentalis principiis omnino elementaribus innixam deteximus.
\par\noindent\hfill [1807] Maii 6.
\par\medskip\begingroup\small
\noindent\hspace*{\parindent}In einem am 8. Mai 1807 begonnenen und am 12. desselben Monats abgeschlossenen Briefe an \textsc{Olbers} (Siehe \textit{Wilhelm Olbers, Sein Leben und seine Werke} II, 1, 1900, S. 357, besonders S. 359, 360) schreibt \textsc{Gauss}:
\par\endgroup
\noindent\hspace*{\parindent}Hierbei alle meine von Mlle. \textsc{Sophie Germain} erhaltenen Briefe. Personalia weiss ich eigentlich weiter keine, als die sich daraus abnehmen lassen. Bloss von einem franz[ösischen] Officier, der im Nov[ember] 1806 hier durchkam, erfuhr ich, dass \textsc{Sophie Germain}, unter der ich aber damals meinen \textsc{Leblanc} noch nicht ahndete, ein in Paris sehr geehrter und bewunderter Name sei. Neulich als ich ihr antwortete und einige Arithmetica \mbox{mit-}
\clearpage\SourceLeaf{596}\ArithmeticHead{567}{R}{1807 FEBR. 22—1808 IUN. 12.}
\begingroup\renewcommand{\thefootnote}{[*)]}
\noindent theilte\footnotemark\ wurde ich dadurch veranlasst, wieder eine Untersuchung vorzunehmen, und gleich zwei Tage nachher gelang mir eine äusserst angenehme neue Entdeckung. Es ist ein neuer, sehr zierlicher und kurzer Beweis des Fundamentalsatzes art. 1[31. der \textit{Disquisitiones arithmeticae}], dessen erster, sehr mühsamer (obwol im Grunde auch einfacher, aber langes Detail erfordernder) Beweis mich über ein Jahr gekostet hatte.\enspace .\enspace .\enspace .
\renewcommand{\thefootnote}{[*)}
\footnotetext{Siehe den Brief vom 30. April 1807, oben S. 70.]}
\par\endgroup
\par\begingroup\small\renewcommand{\thefootnote}{**)}
\noindent\hspace*{\parindent}Nach der durch das \textit{Tagebuch}\footnote{Vergl. die Nr. 118 vom Mai 1801, oben S. 560.} gegebenen Zählung handelt es sich bei der vorliegenden Aufzeichnung und in der Briefstelle an \textsc{Olbers} um den \textls[120]{sechsten} Beweis des Reziprozitätsgesetzes der quadratischen Reste. Die zugehörige Abhandlung führt den Titel \textit{Theorematis arithmetici demonstratio nova} (Werke II, S. 1, vergl. auch die \textit{Anzeige} ebenda S. 151); sie wurde der Königl. Gesellschaft der Wissenschaften am 15. Januar 1808 vorgelegt, also früher als die in den Tagebuchaufzeichnungen Nr. 118 und Nr. 123 angezeigte \textit{Summatio quarumdam serierum singularium} (Werke II, S. 9, vorgelegt am 24. August 1808). Daraus erklärt es sich, daß \textsc{Gauss} die \textit{demonstratio nova} in der \textit{Anzeige}, Werke II, S. 153, abweichend von der Zählung des \textit{Tagebuchs}, als \textls[120]{fünften} Beweis bezeichnet. Vergl. auch den Artikel \textbf{20} des \textsc{Bachmann}schen Aufsatzes »Über \textsc{Gauss}' zahlentheoretische Arbeiten« Werke X2, S. 50.
\par\noindent\hfill\textsc{Klein.}
\par\endgroup

\NumberedItem{135}
\noindent\hspace*{\parindent}Theoria divisionis in periodos \textls[120]{tres} (art. 358) ad principia longe simpliciora reducta.
\par\noindent\hfill 1808 Maii 10.
\par\medskip\begingroup\small
\noindent\hspace*{\parindent}Hier sind wohl die Prinzipien gemeint, die in dem Aufsatze \textit{Disquisitionum circa aequationes puras ulterior evolutio}, Werke II, S. 243, dargelegt sind. Die Artikelnummer bezieht sich auf die \textit{Disquisitiones arithmeticae}, Werke I, S. 445.
\par\noindent\hfill\textsc{Klein.\quad Bachmann.}
\par\endgroup

\NumberedItem{136}
\noindent\hspace*{\parindent}Aequationem
\[
X-1=0,
\]
\noindent quae continet omnes radices primitivas aequationis
\[
x^n-1=0,
\]
\noindent in factores cum coefficientibus rationalibus discerpi non posse, demonstr[atum] pro valoribus compositis ipsius \(n\).
\par\noindent\hfill 1808 Iun. 12.
\par\medskip\begingroup\small
\noindent\hspace*{\parindent}Vergl. die Notiz Nr. 40 vom 9. Oktober 1796, in der die Irreduzibilität der Kreisteilungsgleichung für den Fall, wo \(n\) eine Primzahl ist, angezeigt wird. Daß \textsc{Gauss} die Irreduzibilität für den Fall, wo \(n\)
\par\endgroup
\clearpage\SourceLeaf{597}\ArithmeticHead{568}{L}{TAGEBUCH MIT ERLÄUTERUNGEN.}
\par\begingroup\small
\noindent eine Primzahlpotenz ist, schon bei Abfassung der \textit{Disquisitiones arithmeticae} gekannt haben dürfte, ist in der Bemerkung zu der Nr. 116, oben S. 559, auseinandergesetzt worden. Von dem Beweise der Irreduzibilität im allgemeinen Fall enthält der Nachlaß nur das oben S. 116 abgedruckte Bruchstück, aus dem man aber nicht auf das Beweisverfahren schließen kann, das \textsc{Gauss} im Auge hatte. Auch eine Reihe von Einzelheiten in diesem Bruchstück bleiben unklar, so vor allem die letzten Zeilen von S. 117; vergl. auch die Bemerkung zu der Nr. 38. Im Text der obigen Tagebuchnotiz Nr. 136 dürfte, wie \textsc{Dedekind} beim ersten Abdruck bemerkt hat, die Form \(X-1=0\) der Gleichung ein Schreibfehler für \(X=0\) sein; diese Vermutung wird durch das S. 116 abgedruckte Bruchstück bestätigt. Über die verschiedenen Arten von Irreduzibilitätsbeweisen für die Kreisteilungsgleichung sehe man M. \textsc{Ruthinger}, \textit{Die Irreducibilitätsbeweise der Kreisteilungsgleichung}, Straßburger Dissertation 1907.
\par\noindent\hfill\textsc{Loewy.}
\par\endgroup

\NumberedItem{137}
\noindent\hspace*{\parindent}Theoriam formarum cubicarum, solutionem aequ[ationis]
\[
x^3+ny^3+nnz^3-3nxyz=1
\]
\noindent aggressus sum.
\par\noindent\hfill [1808] Dec. 23.
\par\medskip\begingroup\small
\noindent\hspace*{\parindent}Vergl. die Ausführungen von R. \textsc{Fricke}, Werke VIII, S. 24—26 und eine Bemerkung von \textsc{Schering}, Werke II, S. 398. Die dort zitierte, Werke II, S. 243 abgedruckte Abhandlung \textit{Disquisitionum circa aequationes puras ulterior evolutio} stammt (siehe die Bemerkung von \textsc{Dedekind} a. a. O. S. 265 und die Tagebuchnotiz Nr. 135) in der Tat aus dem Jahre 1808.
\par\noindent\hspace*{\parindent}Vergl. den Artikel \textbf{25} des \textsc{Bachmann}schen Aufsatzes Werke X2, S. 60.
\par\noindent\hfill\textsc{Klein.}
\par\endgroup

\NumberedItem{138}
\noindent\hspace*{\parindent}Theorema de residuo cubico 3 per methodum specialem elegantem demonstratum per considerat[iones] valorum \(\frac{x+1}{x}\), ubi terni semper habent \(a\), \(a\varepsilon\), \(a\varepsilon\varepsilon\) exceptis duobus, qui dant \(\varepsilon\), \(\varepsilon\varepsilon\), hi vero sunt
\[
\frac{1}{\varepsilon-1}=\frac{\varepsilon\varepsilon-1}{3},\quad
\frac{1}{\varepsilon\varepsilon-1}=\frac{\varepsilon-1}{3}
\]
\noindent adeoque productum \(\equiv\frac{1}{3}\).
\par\noindent\hfill 1809 Ian. 6.
\par\medskip\begingroup\small
\noindent\hspace*{\parindent}Hier ist \(\varepsilon\) nicht als dritte Wurzel aus 1, sondern als rationale Wurzel der Kongruenz
\[
\varepsilon^2+\varepsilon+1\equiv0\ \text{(mod.\ }p\text{)}
\]
\noindent aufzufassen, wo \(p\) eine natürliche Primzahl von der Form \(3n+1\) bedeutet.
\par\noindent\hfill\textsc{Dedekind.}
\par\endgroup
